\subsection{A consumption-top-up interpretation}\label{sec:topup}
An absolute utility tolerance is not an economic unit. We therefore define a compensation experiment before reporting it. Keep the policy's preference, wealth, adjustment cost, and stopping paths fixed, and provide an externally financed proportional consumption top-up $\lambda\geq0$. Only the consumption entering instantaneous utility changes, from $c_s$ to $(1+\lambda)c_s$. This is an in-kind welfare comparison; it is not a claim that the extra consumption is financeable under the unchanged budget or the original consumption cap.

\begin{proposition}[A scale-normalized compensation bound]\label{prop:topup}
For a feasible time-control policy with $c\leq1$, $u\in[1.2,2.8]$ until stopping, suppose $V-J^\pi\leq\epsilon$ and $\E\int_0^\tau e^{-.04s}ds\geq A>0$. If $1-1.8\epsilon/A>0$, a top-up no larger than
\begin{equation}
 \overline\lambda=(1-1.8\epsilon/A)^{-1/1.8}-1\label{eq:topup}
\end{equation}
is sufficient to compensate the certified policy loss in the specified welfare experiment.
\end{proposition}
For the fresh central-state price library, the rigorous original stopping calculation supplies
\[
 A\geq\frac{1-e^{-.04}}{.04}\left\{1-2\exp\!\left(-\frac{.58^2}{2(.05)^2}\right)\right\}.
\]
The resulting uniform top-up bound is \RFreshTopupPercent\ percent. It is not assigned to a neural policy. Under a common positive rescaling of the entire welfare criterion, both the loss and the marginal utility gain rescale together, leaving this compensation experiment unchanged. This does not supply an empirical calibration or a wealth-equivalent claim. The exact definition prevents a numerical threshold from being presented as economically meaningful merely because it is slightly below $.01$.
